\documentclass[aps,prl,twocolumn]{revtex4-2}
\usepackage{amsmath,amssymb,booktabs}
\begin{document}

\title{Supplemental Material:\\
Information-Capacity Bound on Non-Markovian Backflow}

\maketitle

\section{Full BLP Derivation}

\subsection{Setup and notation}

$N_S$ system qubits, $N_E$ environment qubits. Initial states $\rho_S^{(1)}(0)$, $\rho_S^{(2)}(0)$ are diagonal in the pointer basis. Define per-qubit populations $q_j^{(m)}(t) \equiv \langle 0|\rho_{S,j}^{(m)}(t)|0\rangle$ for system qubit $j$ and initial state $m \in \{1,2\}$. The trace distance decomposes additively (Assumption A3):
\begin{equation}
D(t) = \sum_{j=1}^{N_S} d_j(t), \qquad d_j(t) \equiv |q_j^{(1)}(t) - q_j^{(2)}(t)|.
\end{equation}

\subsection{Forward transfer does not change D(t)}

Under CNOT$_{S\to E}$, the system qubit is the control. Its reduced state is unchanged. Hence $\dot{q}_j^{(m)} = 0$ during forward transfer, and $dD/dt = 0$. Forward transfer contributes zero to $\mathcal{N}_{\rm BLP}$.

\subsection{Backflow changes D(t)}

A backflow event CNOT$_{E_k \to S_j}$ toggles system qubit $j$ when environment qubit $k$ is in $|1\rangle$. The post-toggle population is:
\begin{equation}
q_j^{(m),{\rm after}} = q_j^{(m)}(1 - P_k^{(m)}) + (1 - q_j^{(m)}) P_k^{(m)},
\end{equation}
where $P_k^{(m)} = {\rm Prob}(E_k = |1\rangle \mid \rho^{(m)})$ is the probability that the controlling environment qubit is in $|1\rangle$.

The per-qubit distance change is $\Delta d_j = |q_j^{(1),{\rm after}} - q_j^{(2),{\rm after}}| - d_j$. The positive part contributes to $\mathcal{N}_{\rm BLP}$. We bound:
\begin{equation}
\Delta d_j^+ \leq |P_k^{(1)} - P_k^{(2)}| \cdot |1 - 2\bar{q}_j|,
\end{equation}
where $\bar{q}_j \equiv (q_j^{(1)} + q_j^{(2)})/2$.

\subsection{Encoding relation}

Under Assumptions A6--A7, the environment qubit's $|1\rangle$-probability after forward encoding is:
\begin{equation}
P_k^{(m)} = q_E + q_j^{(m)} - 2q_E q_j^{(m)}.
\end{equation}
The difference between the two initial states is:
\begin{equation}
P_k^{(1)} - P_k^{(2)} = (q_j^{(1)} - q_j^{(2)})(1 - 2q_E) = \pm d_j \cdot (1 - 2q_E).
\end{equation}

\subsection{Amplification factor and per-event bound}

Combining:
\begin{equation}
\Delta d_j^+ \leq d_j \cdot |1 - 2q_E| \cdot |1 - 2\bar{q}_j|.
\end{equation}
Maximizing over $\bar{q}_j$ subject to the constraint $d_j = 2\bar{q}_j - 2q_j^{(2)}$ (for $q_j^{(1)} > q_j^{(2)}$) yields:
\begin{equation}
\max \Delta d_j^+ = \frac{1}{2}(1 - 2q_E)^2.
\end{equation}
This maximum is achieved at $\bar{q}_j = 3/4$ (for $q_E=0$) and at $\bar{q}_j = 1/2$ (for $q_E=1/2$, giving zero amplification).

\subsection{Integrated BLP bound}

Summing over $N_{\rm back}$ independent backflow events, each affecting one system qubit:
\begin{equation}
\mathcal{N}_{\rm BLP} = \sum_{i=1}^{N_{\rm back}} \Delta D_i^+ \leq N_{\rm back} \cdot \frac{1}{2}(1 - 2q_E)^2.
\end{equation}
Substituting $N_{\rm back} = P_{\rm reflux} \cdot C_F = P_{\rm reflux} \cdot N_E \cdot q_E$ and applying the S1 bound $P_{\rm reflux} \leq q_S/q_E$:
\begin{equation}
\boxed{\mathcal{N}_{\rm BLP} \leq \frac{1}{2} \cdot P_{\rm reflux} \cdot N_E \cdot q_E \cdot (1 - q_S) \cdot (1 - 2q_E)^2 \leq \frac{N_E}{2} \cdot q_S \cdot (1 - q_S) \cdot (1 - 2q_E)^2}.
\end{equation}

\subsection{Relaxation of Assumptions A6--A7}

If Assumption A6 is relaxed (environment qubits undergo multiple forward CNOTs), the encoding degrades by a factor $\exp(-\gamma N_{\rm forward}/N_E)$. The bound acquires a multiplicative correction $\eta \in (0,1]$. The qualitative structure $\mathcal{N}_{\rm BLP} \propto q_S(1-q_S)$ is preserved.

If Assumption A7 is relaxed (forward encoding is inefficient), replace $q_E \to \eta_E q_E$ where $\eta_E$ is the encoding efficiency. The bound weakens by a constant factor but remains structurally identical.

\section{Quantum Darwinism and Irreversibility}

\subsection{Kac recurrence time estimate}

For $N_{\rm red}$ environment qubits redundantly encoding the $|1\rangle$ state of a determined system qubit, the joint state of the $N_{\rm red}$ records is approximately $|1\rangle^{\otimes N_{\rm red}}$. This is one of $2^{N_{\rm red}}$ possible configurations. The Kac recurrence time---the expected time for unitary evolution to return to this state---is $\tau_{\rm Kac} \sim 2^{N_{\rm red}} \cdot \tau_{\rm micro}$, where $\tau_{\rm micro} \sim \hbar/\Delta E$ is the characteristic microscopic timescale. For $\tau_{\rm micro} \sim 10^{-15}\,{\rm s}$ and $N_{\rm red} \gtrsim 100$, $\tau_{\rm Kac} \gtrsim 10^{30}\,{\rm yr}$, exceeding the age of the universe by $\sim 20$ orders of magnitude. On all physically relevant timescales, the $|0\rangle \to |1\rangle$ transition is effectively irreversible.

\subsection{Relation to no-cloning}

The irreversibility is consistent with the no-cloning theorem: the forward CNOT copies \emph{classical} information (a definite bit value in the pointer basis), not an arbitrary quantum state. Classical information can be copied; quantum information cannot. Since decoherence has already diagonalized the state in the pointer basis before forward transfer, the information encoded is classical, and the CNOT operation is a classical COPY, not a quantum cloning attempt.

\section{Kraus Operator Construction for Forward Transfer}

The forward transfer quantum instrument $\mathcal{E}_{FT}$ has Kraus operators:
\begin{align}
K_0 &= |0\rangle\langle 0|_S \otimes I_E, \\
K_1 &= |1\rangle\langle 1|_S \otimes |1\rangle\langle 0|_E.
\end{align}
These satisfy $\sum_i K_i^\dagger K_i = I_{SE}$ (completeness). The operation is:
\begin{equation}
\mathcal{E}_{FT}(\rho) = K_0 \rho K_0^\dagger + K_1 \rho K_1^\dagger.
\end{equation}
When $S = |0\rangle$, only $K_0$ acts (no transfer). When $S = |1\rangle$ and $E = |0\rangle$, $K_1$ toggles $E$ to $|1\rangle$. When $E$ is already $|1\rangle$, $K_1$ annihilates the state (the operation is not defined for this input), reflecting the physical fact that a determined environment qubit cannot be re-determined.

\section{Operational Definition of $q_S$ and $q_E$}

$q_S$ and $q_E$ are measured via projective measurement in the pointer basis. The pointer basis is identified experimentally as the basis in which the system's reduced state remains diagonal under the system-environment interaction---equivalently, the basis that minimizes the decoherence rate. For each qubit, measure $\sigma_z$ repeatedly. The fraction of $+1$ outcomes (corresponding to $|0\rangle$) over many runs gives $\langle 0|\rho|0\rangle$. The per-qubit average over all system qubits gives $q_S$.

\begin{thebibliography}{5}
\bibitem{Breuer:2009} H.-P.~Breuer, E.-M.~Laine, and J.~Piilo, Phys.\ Rev.\ Lett.\ \textbf{103}, 210401 (2009).
\bibitem{Zurek:2003} W.~H.~Zurek, Rev.\ Mod.\ Phys.\ \textbf{75}, 715 (2003).
\bibitem{Zurek:2009} W.~H.~Zurek, Nature Phys.\ \textbf{5}, 181 (2009).
\end{thebibliography}

\end{document}
